\documentclass[10pt,a4paper]{amsart}
\usepackage{amsmath,amssymb,amsthm,graphicx}
\usepackage{verbatim}
\newenvironment{reference}{\comment}{\endcomment}
%\newenvironment{reference}{}{}
\newenvironment{slogan}{\comment}{\endcomment}
\newenvironment{history}{\comment}{\endcomment}

% For commutative diagrams we use Xy-pic
\usepackage[all]{xy}

% We use 2cell for 2-commutative diagrams.
\xyoption{2cell}
% Theorem environments.
%
\theoremstyle{plain}
\newtheorem{theorem}[subsection]{Theorem}
\newtheorem{proposition}[subsection]{Proposition}
\newtheorem{lemma}[subsection]{Lemma}

\theoremstyle{definition}
\newtheorem{definition}[subsection]{Definition}
\newtheorem{example}[subsection]{Example}
\newtheorem{exercise}[subsection]{Exercise}
\newtheorem{situation}[subsection]{Situation}

\theoremstyle{remark}
\newtheorem{remark}[subsection]{Remark}
\newtheorem{remarks}[subsection]{Remarks}

\numberwithin{equation}{subsection}

% Macros
%
\def\lim{\mathop{\mathrm{lim}}\nolimits}
\def\colim{\mathop{\mathrm{colim}}\nolimits}
\def\Spec{\mathop{\mathrm{Spec}}}
\def\Hom{\mathop{\mathrm{Hom}}\nolimits}
\def\Ext{\mathop{\mathrm{Ext}}\nolimits}
\def\SheafHom{\mathop{\mathcal{H}\!\mathit{om}}\nolimits}
\def\SheafExt{\mathop{\mathcal{E}\!\mathit{xt}}\nolimits}
\def\Sch{\mathit{Sch}}
\def\Mor{\mathop{\mathrm{Mor}}\nolimits}
\def\Ob{\mathop{\mathrm{Ob}}\nolimits}
\def\Sh{\mathop{\mathit{Sh}}\nolimits}
\def\NL{\mathop{N\!L}\nolimits}
\def\CH{\mathop{\mathrm{CH}}\nolimits}
\def\proetale{{pro\text{-}\acute{e}tale}}
\def\etale{{\acute{e}tale}}
\def\QCoh{\mathit{QCoh}}
\def\Ker{\mathop{\mathrm{Ker}}}
\def\Im{\mathop{\mathrm{Im}}}
\def\Coker{\mathop{\mathrm{Coker}}}
\def\Coim{\mathop{\mathrm{Coim}}}

% Boxtimes
%
\DeclareMathSymbol{\boxtimes}{\mathbin}{AMSa}{"02}

%
% Macros for moduli stacks/spaces
%
\def\QCohstack{\mathcal{QC}\!\mathit{oh}}
\def\Cohstack{\mathcal{C}\!\mathit{oh}}
\def\Spacesstack{\mathcal{S}\!\mathit{paces}}
\def\Quotfunctor{\mathrm{Quot}}
\def\Hilbfunctor{\mathrm{Hilb}}
\def\Curvesstack{\mathcal{C}\!\mathit{urves}}
\def\Polarizedstack{\mathcal{P}\!\mathit{olarized}}
\def\Complexesstack{\mathcal{C}\!\mathit{omplexes}}
% \Pic is the operator that assigns to X its picard group, usage \Pic(X)
% \Picardstack_{X/B} denotes the Picard stack of X over B
% \Picardfunctor_{X/B} denotes the Picard functor of X over B
\def\Pic{\mathop{\mathrm{Pic}}\nolimits}
\def\Picardstack{\mathcal{P}\!\mathit{ic}}
\def\Picardfunctor{\mathrm{Pic}}
\def\Deformationcategory{\mathcal{D}\!\mathit{ef}}
\usepackage{cleveref}
\crefname{originallemma}{original lemma}{original lemmas}
\Crefname{originallemma}{Original lemma}{Original lemmas}
\makeatletter
\newlabel{lemma-NAK}{{00DV}{source}{}{sourcepointer-1}{}}
\newlabel{lemma-NAK@cref}{{[originallemma][0][]00DV}{[1][0][]source}{}{}{}}
\makeatother
\crefname{originallemma}{original lemma}{original lemmas}
\Crefname{originallemma}{Original lemma}{Original lemmas}
\makeatletter
\newlabel{lemma-characteristic-zero}{{00TW}{source}{}{sourcepointer-2}{}}
\newlabel{lemma-characteristic-zero@cref}{{[originallemma][0][]00TW}{[1][0][]source}{}{}{}}
\makeatother
\crefname{originallemma}{original lemma}{original lemmas}
\Crefname{originallemma}{Original lemma}{Original lemmas}
\makeatletter
\newlabel{lemma-computation-differential}{{00TU}{source}{}{sourcepointer-3}{}}
\newlabel{lemma-computation-differential@cref}{{[originallemma][0][]00TU}{[1][0][]source}{}{}{}}
\makeatother
\crefname{originallemma}{original lemma}{original lemmas}
\Crefname{originallemma}{Original lemma}{Original lemmas}
\makeatletter
\newlabel{lemma-differential-seq}{{00RU}{source}{}{sourcepointer-4}{}}
\newlabel{lemma-differential-seq@cref}{{[originallemma][0][]00RU}{[1][0][]source}{}{}{}}
\makeatother
\crefname{originallemma}{original lemma}{original lemmas}
\Crefname{originallemma}{Original lemma}{Original lemmas}
\makeatletter
\newlabel{lemma-regular-mod-x}{{00NU}{source}{}{sourcepointer-5}{}}
\newlabel{lemma-regular-mod-x@cref}{{[originallemma][0][]00NU}{[1][0][]source}{}{}{}}
\makeatother
\crefname{originallemma}{original lemma}{original lemmas}
\Crefname{originallemma}{Original lemma}{Original lemmas}
\makeatletter
\newlabel{lemma-separable-smooth}{{00TV}{source}{}{sourcepointer-6}{}}
\newlabel{lemma-separable-smooth@cref}{{[originallemma][0][]00TV}{[1][0][]source}{}{}{}}
\makeatother
\title{In characteristic zero smoothness at a prime, freeness of the differentials and regularity agree}
\author{The Stacks Project Authors}
\date{stacks-project:algebra.tex:local-snapshot-20261012; GNU FDL 1.2 or later}
\begin{document}
\maketitle
\noindent\textit{Editorial note.} All delivered mathematical text is the author's own, copied from the pinned Stacks Project chapter file algebra.tex; no statement or proof step was rewritten, reordered or removed.\par\medskip
\noindent\textit{Editorial note.} Preamble provenance, as required for this source: the frozen snapshot stages the chapter's own preamble file beside it (preamble.tex, 217 lines, 8116 bytes, the same shared Stacks-wide preamble the published chow and etale-cohomology cards use); the first line of the pinned chapter file is \textbackslash{}input\{preamble\}, so no preamble text lives inside the chapter itself. The card takes that file's lines 8-18 (the verbatim package and the three \textbackslash{}comment-based environments reference/slogan/history at lines 9-12, section 13.2, together with the Xy-pic setup) and its lines 154 to the runtime end of the file (the native theorem declarations and mathematical macros). Not one character was changed, no package was added, no macro was defined, and the excerpt serves only as preamble, never as a span coordinate.\par\medskip
\noindent\textit{Editorial note.} The same-chapter pointers are located at their \textbackslash{}label lines inside the pinned chapter file; the cross-chapter pointers are located in the project's own tag list staged in the snapshot (sources/stacks-chow/tags.txt), at the verbatim lines carrying their canonical tags.\par\medskip
\noindent\textit{Editorial note.} The counted proof refers to same-chapter statements by label; they are recorded as navigation stubs pointed at their label lines.\par\medskip
\noindent\textit{Editorial note.} The counted spans contain no cite, so the card carries no bibliography.\par\medskip
\noindent\textit{Allowed standard prerequisites (editorial).}\begin{itemize}
\item Separability of field extensions in characteristic zero. Used at the start of the counted proof (lines 39052-39053). Gives the equivalence between smoothness and regularity through the separable case.
\item Regularity inherited after dividing by a nonzerodivisor. Invoked at the end of the counted proof (lines 39081-39083). Lifts regularity from a quotient along a nonzerodivisor.
\end{itemize}\medskip
\ifdefined\hypertarget\hypertarget{sourcepointer-1}{}\fi
\noindent\textit{Source reference.} Original lemma 00DV: sources/stacks-algebra/pinned-original-source.tex, source lines 3833--3833. Navigation pointer for \textbackslash{}ref\{lemma-NAK\} used inside the counted span (line 39072); the argument that consumes it is retained in the span.\par\medskip
\ifdefined\hypertarget\hypertarget{sourcepointer-2}{}\fi
\noindent\textit{Source reference.} Original lemma 00TW: sources/stacks-algebra/pinned-original-source.tex, source lines 39001--39001. Navigation pointer for \textbackslash{}ref\{lemma-characteristic-zero\} used inside the counted span (line 39074); the argument that consumes it is retained in the span.\par\medskip
\ifdefined\hypertarget\hypertarget{sourcepointer-3}{}\fi
\noindent\textit{Source reference.} Original lemma 00TU: sources/stacks-algebra/pinned-original-source.tex, source lines 38904--38904. Navigation pointer for \textbackslash{}ref\{lemma-computation-differential\} used inside the counted span (line 39068); the argument that consumes it is retained in the span.\par\medskip
\ifdefined\hypertarget\hypertarget{sourcepointer-4}{}\fi
\noindent\textit{Source reference.} Original lemma 00RU: sources/stacks-algebra/pinned-original-source.tex, source lines 34320--34320. Navigation pointer for \textbackslash{}ref\{lemma-differential-seq\} used inside the counted span (line 39076); the argument that consumes it is retained in the span.\par\medskip
\ifdefined\hypertarget\hypertarget{sourcepointer-5}{}\fi
\noindent\textit{Source reference.} Original lemma 00NU: sources/stacks-algebra/pinned-original-source.tex, source lines 26033--26033. Navigation pointer for \textbackslash{}ref\{lemma-regular-mod-x\} used inside the counted span (line 39083); the argument that consumes it is retained in the span.\par\medskip
\ifdefined\hypertarget\hypertarget{sourcepointer-6}{}\fi
\noindent\textit{Source reference.} Original lemma 00TV: sources/stacks-algebra/pinned-original-source.tex, source lines 38957--38957. Navigation pointer for \textbackslash{}ref\{lemma-separable-smooth\} used inside the counted span (line 39053); the argument that consumes it is retained in the span.\par\medskip
\begin{lemma}
\label{lemma-characteristic-zero-local-smooth}
Let $k$ be a field of characteristic $0$.
Let $S$ be a finite type $k$-algebra.
Let $\mathfrak q \subset S$ be a prime.
The following are equivalent:
\begin{enumerate}
\item The algebra $S$ is smooth at $\mathfrak q$ over $k$.
\item The $S_{\mathfrak q}$-module $\Omega_{S/k, \mathfrak q}$
is (finite) free.
\item The ring $S_{\mathfrak q}$ is regular.
\end{enumerate}
\end{lemma}

\begin{proof}
In characteristic zero any field extension is separable and hence the
equivalence of (1) and (3) follows from Lemma \ref{lemma-separable-smooth}.
Also (1) implies (2) by definition of smooth algebras.
Assume that $\Omega_{S/k, \mathfrak q}$ is free over $S_{\mathfrak q}$.
We are going to use the notation and observations made in the
proof of Lemma \ref{lemma-separable-smooth}. So $R = S_{\mathfrak q}$
with maximal ideal $\mathfrak m$ and residue field $\kappa$.
Our goal is to prove $R$ is regular.

\medskip\noindent
If $\mathfrak m/\mathfrak m^2 = 0$, then $\mathfrak m = 0$
and $R \cong \kappa$. Hence $R$ is regular and we win.

\medskip\noindent
If $\mathfrak m/ \mathfrak m^2 \not = 0$, then choose any
$f \in \mathfrak m$ whose image in $\mathfrak m/ \mathfrak m^2$
is not zero. By Lemma \ref{lemma-computation-differential}
we see that $\text{d}f$ has nonzero image in
$\Omega_{R/k}/\mathfrak m\Omega_{R/k}$. By assumption
$\Omega_{R/k} = \Omega_{S/k, \mathfrak q}$ is finite free and
hence by Nakayama's Lemma \ref{lemma-NAK} we see that
$\text{d}f$ generates a direct summand. We apply
Lemma \ref{lemma-characteristic-zero}
to deduce that $f$ is a nonzerodivisor in $R$.
Furthermore, by Lemma \ref{lemma-differential-seq} we get an exact sequence
$$
(f)/(f^2) \to \Omega_{R/k} \otimes_R R/fR \to \Omega_{(R/fR)/k} \to 0
$$
This implies that $\Omega_{(R/fR)/k}$ is finite free as well.
Hence by induction we see that $R/fR$ is a regular local ring.
Since $f \in \mathfrak m$ was a nonzerodivisor we
conclude that $R$ is regular, see Lemma \ref{lemma-regular-mod-x}.
\end{proof}

\end{document}
